% ============================================================
% CAPITOLO 12 - LE CADENZE
% ============================================================
\chapter{Le cadenze}
\label{chap:cadenze}

Se gli accordi sono le parole del linguaggio musicale, le \textbf{cadenze} ne sono la punteggiatura.

\begin{definizione}[La cadenza]
	\label{def:cadenza}
	Una \textbf{cadenza} è una formula melodico-armonica fatta di due o più accordi che crea un punto di articolazione, pausa o conclusione nel discorso musicale.
\end{definizione}

Le cadenze delimitano le frasi e strutturano il flusso temporale. Alcune danno un senso di \textbf{conclusione} e riposo, altre di \textbf{sospensione} e attesa, e così il compositore guida le aspettative dell'ascoltatore. La loro efficacia dipende dalla successione degli accordi, dalla condotta delle voci, dal ritmo armonico e dalla posizione metrica.

\section{Cadenze conclusive}
\label{sec:cadenze_conclusive}

Le cadenze conclusive producono chiusura e riposo e servono a terminare le frasi in modo definitivo.

\begin{definizione}[Cadenza perfetta]
	\label{def:cadenza_perfetta}
	La \textbf{cadenza perfetta} (o autentica perfetta) è la formula conclusiva più forte del sistema tonale. È la successione \textbf{V - I}, entrambi in stato fondamentale, con la tonica al soprano nell'accordo finale.
\end{definizione}

La dominante, con la sensibile che spinge verso l'alto, risolve nel centro di gravità della tonica e porta il discorso a una quiete completa.

\begin{esempio}[Cadenza perfetta in Do maggiore]
	\label{ex:cadenza_perfetta}
	\musica{es_cap06_01}
\end{esempio}

\begin{definizione}[Cadenza imperfetta]
	\label{def:cadenza_imperfetta}
	La \textbf{cadenza imperfetta} (o autentica imperfetta) è una cadenza V-I (o \sigla{V}{7}{}-I) in cui si verifica almeno una di queste condizioni:
	\begin{itemize}
		\item il V non è in stato fondamentale (per esempio \sigla{V}{6}{} o \sigla{V}{6}{5});
		\item il I non è in stato fondamentale (per esempio \sigla{I}{6}{});
		\item il soprano dell'accordo di arrivo non ha la tonica ma la terza o la quinta.
	\end{itemize}
	Dà una chiusura meno definitiva della perfetta ed è usata spesso a metà periodo (Capitolo \ref{chap:struttura_frase}).
\end{definizione}

\begin{esempio}[Cadenza imperfetta]
	Due varianti: soprano sulla terza (V-I) e \sigla{V}{6}{} al posto di V (\sigla{V}{6}{}-I).
	\musica{es_cap06_02}
	\musica{es_cap06_03}
\end{esempio}

\begin{definizione}[Cadenza plagale]
	\label{def:cadenza_plagale}
	La \textbf{cadenza plagale} è la successione \textbf{IV - I}. Dà una chiusura più debole e statica della perfetta. Si usa spesso come formula aggiuntiva dopo una cadenza perfetta ed è nota come la cadenza dell'«Amen» nella musica liturgica.
\end{definizione}

\begin{esempio}[Cadenza plagale in Do maggiore]
	\label{ex:cadenza_plagale}
	\musica{es_cap06_04}
\end{esempio}

\begin{regola}[Gerarchia delle cadenze]
	\label{reg:gerarchia_cadenze}
	La forza conclusiva decresce in quest'ordine:
	\begin{center}
		\textbf{C. perfetta > C. imperfetta > C. plagale > C. sospesa > C. d'inganno}
	\end{center}
	La gerarchia guida la scelta della cadenza in base al posto che occupa nella struttura formale (Capitolo \ref{chap:struttura_frase}).
\end{regola}

\section{Cadenze sospensive}
\label{sec:cadenze_sospensive}

Le cadenze sospensive creano interruzione e attesa e lasciano il discorso «aperto». Sono l'equivalente musicale di una virgola o di un punto interrogativo.

\begin{definizione}[Cadenza sospesa]
	\label{def:cadenza_sospesa}
	La \textbf{cadenza sospesa} (o semicadenza) è una formula armonica che termina sulla \textbf{dominante (V)}. L'accordo precedente può essere la tonica (I), la sottodominante (IV), il II grado o altri. Crea l'attesa della risoluzione sulla tonica, che viene rimandata.
\end{definizione}

\begin{esempio}[Cadenza sospesa (IV - V)]
	\label{ex:cadenza_sospesa}
	\musica{es_cap06_05}
\end{esempio}

\begin{definizione}[Cadenza d'inganno]
	\label{def:cadenza_inganno}
	La \textbf{cadenza d'inganno} è la successione \textbf{V - VI}. L'ascoltatore si aspetta la tonica dopo la dominante e trova invece il VI grado, che ha due note in comune con la tonica ma ne sposta il centro: nasce una sorpresa che prolunga la frase.
\end{definizione}

Nella cadenza d'inganno la sensibile sale comunque alla tonica e il basso sale di grado congiunto. Per evitare quinte parallele si raddoppia di norma la terza dell'accordo del VI grado.

\begin{esempio}[Cadenza d'inganno in Do maggiore]
	\label{ex:cadenza_inganno}
	\musica{es_cap06_06}
\end{esempio}
